\scanpage{90}
\margin{89}
qu'à remplacer $\Gamma$ par un sous-groupe ouvert convenable $\Gamma'$\ldots)

En $=$ temps, on entrevoit des conditions nécessaires draconiennes pour qu'un autom. \inserted{$\Theta$} du groupoïde $\Pi_1(U_{0,3};I_0)$, \unsure{laiss.$^t$|laiss.} \unsure{fixe|fixes} les sommets i.e. les pts de $I_0$, provienne d'un $u\in\Gamma$: il faut que pour $\forall$ revêtement \deleted{de $X$} étale $C$ de $U_{0,3}$, défini sur $K$, \inserted{prenant $I=f^{-1}(I_0)\subset C_K$ et \unsure{projetant|proj.} $I\to I_0$, \unsure{surj.}} \deleted{telle que} \deleted{l'action de} \deleted{\unsure{l'opération de}} \inserted{\unsure{que l'autom.}} $\Theta\deleted{_{\unsure{C_0}}}$ \deleted{de} $\Pi_1(U_{0,3},I_0)$ \deleted{provienne d'une action de} \inserted{soit ``compatible'' avec un autom.} $\Theta_C$ de $\Pi_1(C,I)$ \inserted{qui soit l'identité} \deleted{\unsure{qui soit ponctuel|qui soit ponctuellement}} \inserted{sur les sommets $I$} \deleted{par l'identi} (cet $\Theta_C$ sera alors unique). On trouve encore une condition, $=$ si $I-I_0$ n'est pas surjectif, et \deleted{en} \unsure{supposons|supposant} $=$ que $I_0\subset C_{0K}$ pour $\forall$ rev. étale de $U_{0,3}$, en exigeant qu'il existe une puissance $\Theta^n$ de $\Theta$ \inserted{($n\geq1$)} satisfaisant aux conditions précédentes, pour $I=f^{-1}(I_0)$.
\margin{(NB $C_0=U_{0,3}=P^1_{\mathbb C}\setminus\{0,1,\infty\}$)}
\margin{\unsure{Regarder} des \unsure{familles} \unsure{pour} $\unsure{C_0}$ \illegible{2} \deleted{\unsure{telles}} \unsure{conditions} \illegible{2}, quand on augmente $I_0$\ldots}
